\clearpage
\section{تشریح موضوع و بیان مسئله}
\subsection{بیان مسئله}

یک تصویر فوندوس همه‌ی اطلاعات لازم برای تصمیم درمانی را در اختیار نمی‌گذارد، اما می‌تواند برای رتبه‌بندی موارد مشکوک به کار رود. آنچه سنجیده می‌شود فقط توان رتبه‌بندی مدل‌ها در برابر برچسب ثبت‌شده در داده است، نه کاربرد بالینی، نه ایمنی و نه سودمندی در عمل.

پرسش نخست این است که آیا اندازه‌گیری صریح عروق، \lr{Plus}\enfoot{گشادشدگی و پیچ‌خوردگی پیش‌رونده‌ی عروق قطب خلفی؛ معادل انگلیسی: \lr{Plus}} را از دو وضعیت دیگر جدا می‌کند و پرسش دوم این است که آیا افزودن همین اندازه‌گیری به یک بازنمایی تصویری آموخته‌شده رتبه‌بندی را بهتر می‌کند. پاسخ این دو لزوماً یکی نیست، پس اندازه‌گیری عروقی در دو سطح آزموده شد، پنج نشانگر اسکالر و بازنمایی فضایی نقشه‌ی عروق، و مقایسه‌ها روی تصاویر یکسان، یک تقسیم و یک معیار انجام گرفت.

\subsection{چالش‌های کلیدی}

نخستین چالش وابستگی نمونه‌هاست. تصاویر یک نوزاد شبیه هم‌اند و اگر میان آموزش و آزمون پخش شوند، مدل هویت چشم یا دستگاه را به‌جای نشانه‌ی بیماری به خاطر می‌سپارد. تقسیم این پروژه گروه‌محور است و هیچ گروهی در دو بخش نیست (جدول \autoref{tab:cohort}). تعریف گروه اما در همه‌ی منابع یکسان نیست؛ در \lr{Plus} و \lr{FARFUM-RoP} در سطح بیمار و در فارابی در سطح معاینه، چون شناسه‌ی بیمار در آن منبع وجود نداشت. پس در فارابی، نبود هم‌پوشانی معاینه نبود هم‌پوشانی بیمار را اثبات نمی‌کند.

چالش دوم وابستگی به منبع تصویربرداری است. سه منبع ترکیب کلاس‌های بسیار متفاوتی دارند و منبع \lr{Plus} هیچ مورد \lr{Pre-Plus} ندارد؛ پس فقط معیار دودویی محدود \lr{Normal} در برابر \lr{Plus} با همان برچسب صریح گزارش می‌شود. این وابستگی اندازه‌گیری شد؛ جابه‌جایی بازنمایی‌ها میان هر منبع کنارگذاشته و منابع آموزش بدون استفاده از برچسب هدف محاسبه شد و \lr{Plus} جابه‌جاشده‌ترین منبع بود (جدول \autoref{tab:domainshift}).

چالش سوم به کیفیت اندازه‌گیری نشانگرها مربوط است. نشانگرها از ماسک عروق می‌آیند و خطای جداسازی به خطای اندازه‌گیری تبدیل می‌شود. ماسک مرجع متخصص برای داده‌ی این پروژه در دسترس نبود، پس صحت جداسازی در برابر مرز انسانی سنجیده نشد و کنترل‌ها فقط سلامت فنی خروجی را تأیید می‌کنند؛ یک نشانگر ثابت هم در سه مرکز عدد یکسانی به دست نمی‌دهد.

چالش چهارم انضباط آزمون است. اگر مدل، پیکربندی یا آستانه پس از دیدن آزمون تغییر کند، آن آزمون برآورد مستقلی از عملکرد نمی‌دهد. در این پروژه انتخاب مدل و آستانه فقط با اعتبارسنجی انجام شد و آزمون قفل‌شده برای نتیجه‌ی اصلی به کار رفت؛ تحلیل‌های بعدی روی همان تقسیم، بازتحلیل اکتشافی ثانویه‌اند.

\subsection{رویکرد این پروژه}

طراحی نهایی روی یک تقسیم استاندارد مشترک بنا شده است که در سطح گروه ساخته شد و هیچ گروه مشترکی میان بخش‌هایش نیست (جدول \autoref{tab:cohort}). همه‌ی مدل‌ها روی همین تقسیم اجرا شدند تا اختلاف معیارها به مدل نسبت داده شود، نه به داده.

سه فرمول‌بندی ویژگی محور مقایسه‌اند؛ بازنمایی تصویری، تصویر به‌همراه بازنمایی فضایی عروق، و تصویر به‌همراه عروق و پنج نشانگر اسکالر. دو مرجع لازم هم اضافه می‌شود؛ نشانگرهای اسکالر به‌تنهایی و نقشه‌ی عروق به‌تنهایی، و نسخه‌ی الحاق زودهنگام نشانگرها برای همان پرسش ترکیب اجرا شد. این مجموعه پرسش «آیا اندازه‌گیری صریح عروقی اطلاعات دارد» را از پرسش «آیا افزودن آن به بازنمایی آموخته‌شده سود دارد» جدا می‌کند. استنباط بر بوت‌استرپ زوجی طبقه‌بندی‌شده استوار است و چون نمونه‌ها زوج‌اند، نتیجه به نابرابری اندازه‌ی کلاس‌ها حساس نیست.

انتقال با سه فولد منبع‌کنارگذاشته سنجیده شد؛ در هر فولد یک منبع کامل از آموزش، پیش‌پردازش و انتخاب مدل کنار گذاشته می‌شود و هیچ ویژگی بدون برچسبی از منبع هدف استفاده نمی‌شود. روی همین بستر، یک آزمایش ناوردایی دامنه‌ی شرطی‌به‌کلاس با ضرایب ثابت و بدون تنظیم روی منبع هدف اجرا شد؛ یک کنترل بی‌طرف در برابر همان شبکه به‌همراه متمایزکننده‌ی دامنه\enfoot{شبکه‌ی عصبی تقابلی دامنه با لایه‌ی معکوس‌کننده‌ی گرادیان؛ معادل انگلیسی: \lr{Domain-Adversarial Neural Network}} و فاصله‌ی \lr{MMD}\enfoot{فاصله‌ی بیشینه‌ی میانگین‌ها میان دو توزیع؛ معادل انگلیسی: \lr{Maximum Mean Discrepancy}} شرطی‌به‌کلاس. قابلیت پیش‌بینی هویت دامنه هم اندازه‌گیری شد و آخرین حلقه‌ی طراحی، تکرار پرسش عروقی با یک ستون فقرات تصویری دوم است.
